% chapters/ch12-relationships-advanced.tex
\chapter{Relationships in Practice: Optionality, Specialization, and Association}

\begin{learningobjectives}
In this chapter you will deepen your understanding of relationships.
You will learn how optionality affects schema structure, how to model specialization (subtypes) without ambiguity, and how to treat association relations as first-class parts of the model when they carry meaning.
\end{learningobjectives}

\section{Optionality is a design decision}
Optionality is often treated as a small detail: “nullable or not”.
In reality, optionality is a statement about what the domain allows.
It affects correctness, lifecycle, and how you validate data.

Consider the sentence: “An appointment has a doctor.”
This implies mandatory participation.
But consider: “A patient may have an insurance profile.”
This implies optional participation.
The difference is not about storage.
It is about whether the system accepts the state “patient exists but no insurance profile exists”.

Optionality decisions should be explicit in the rules.
If optionality is left implicit, the schema will often drift into permissiveness, allowing states that the domain does not allow.

\begin{examplebox}
\textbf{Optionality in words}\\
Mandatory: “Each appointment is with exactly one doctor.”\\
Optional: “A patient may have an insurance profile.”
\end{examplebox}

\subsection{Optionality and foreign key placement}
In one-to-one relationships, optionality guides where the foreign key goes.
If one side can exist without the other, place the foreign key on the optional side and enforce uniqueness there.
This reduces nulls and matches typical creation workflows.

In one-to-many relationships, optionality decides whether the foreign key can be missing.
For example, if an appointment must reference a doctor, \texttt{doctor\_id} is mandatory.
If you make it optional, you allow “unassigned appointments”, which might be correct in some domains and incorrect in others.
The schema should follow the rule, not the other way around.

\section{Association relations are not “just join tables”}
Many-to-many relationships are represented with association relations.
Beginners often treat these as purely technical structures, but they are often conceptually central.
An association relation is not a workaround.
It is the correct representation of a relationship that has its own identity.

When an association has attributes, it becomes clearly meaningful.
In a university domain, enrollment is the association between student and course, and it has attributes such as enrollment date and grade.
In a clinic domain, doctor specialties may include the date the doctor was certified.

Even when an association has no attributes, it still represents meaning: membership, eligibility, assignment, coverage.

\begin{definitionbox}
An \term{association relation} represents a many-to-many relationship.
It becomes conceptually central when the relationship carries attributes or when it represents an event or membership over time.
\end{definitionbox}

\begin{examplebox}
\textbf{Enrollment as a first-class association}\par\smallskip
\[
\begin{aligned}
\rel{Enrollment}(&\underline{student\_id} \rightarrow \rel{Student}, \\
&\underline{course\_id} \rightarrow \rel{Course}, enrolled\_at, grade?)
\end{aligned}
\]
Enrollment is not a technical join; it is a meaningful entity in the domain.
\end{examplebox}

A good teaching move is to ask: “If I delete this association row, what real-world thing disappears?”
If the answer is “an enrollment” or “a certification” or “a membership”, then you are looking at a real concept, not a technical artefact.

\section{Specialization: modelling subtypes}
Many domains contain “kinds of” entities.
A user can be a student or a staff member.
An appointment can be an exam or a consultation.
A payment can be a card payment or a bank transfer.

This is called \term{specialization} or \term{subtyping}.
The key modelling question is: do the subtypes have distinct attributes or constraints that matter?
If not, they can remain as a simple attribute.
If yes, subtyping becomes useful.

\begin{definitionbox}
A \term{specialization} models a supertype entity with one or more subtypes.
Subtypes inherit identity from the supertype and add subtype-specific attributes or constraints.
\end{definitionbox}

\begin{figure}[ht]
\centering
\includegraphics[width=0.75\textwidth]{specialization.pdf}
\caption{A specialization shows a supertype with subtypes linked by an is-a relationship.}
\end{figure}

\subsection{Two constraints: disjointness and completeness}
Subtype models have two structural rules.

Disjointness answers: can an instance belong to more than one subtype?
A person might be both a staff member and a student, depending on the domain.
In other domains, a vehicle is either a car or a motorcycle, but not both.

Completeness answers: must every supertype instance belong to at least one subtype?
Sometimes yes (every payment is exactly one payment type).
Sometimes no (you may create a person record before deciding their role).

These rules must be stated.
If they are not, the subtype model becomes ambiguous.

\begin{examplebox}
\textbf{Subtype constraints in words}\\
Disjoint: “An appointment is either an exam or a consultation, not both.”\\
Overlapping: “A person may be both a student and a staff member.”\\
Complete: “Every payment must have exactly one payment subtype.”\\
Incomplete: “A person may exist without a role assigned yet.”
\end{examplebox}

\section{Implementing subtypes in a relational schema (concept-first)}
There are three common relational patterns for subtypes.
We describe them conceptually without committing to SQL.

\subsection{Pattern 1: Single table with a type attribute}
You store all attributes in one relation and add a type attribute.
This works when subtype-specific attributes are few, optional, and not essential for correctness.

Example:
\[
\begin{aligned}
\rel{Appointment}(&\underline{appointment\_id}, \dots, \\
&type \in \{\text{exam, consultation}\}, exam\_room?, exam\_code?)
\end{aligned}
\]
The risk is that you end up with many nullable columns and weak enforcement of subtype constraints.

\subsection{Pattern 2: Supertype + subtype tables (shared primary key)}
You create one relation for the supertype and one for each subtype.
Subtype relations use the same primary key value as the supertype.
This is clean when subtype attributes matter and you want strong structure.

\begin{examplebox}
\textbf{Supertype + subtype schema}\par\smallskip
\[
\rel{Payment}(\underline{payment\_id}, amount, created\_at)
\]
\[
\rel{CardPayment}(\underline{payment\_id} \rightarrow \rel{Payment}, last4, auth\_code)
\]
\[
\rel{BankTransfer}(\underline{payment\_id} \rightarrow \rel{Payment}, iban, reference)
\]
\end{examplebox}

This pattern represents inheritance clearly.
It also matches the conceptual idea: every card payment is a payment.

\subsection{Pattern 3: Separate tables per subtype (no explicit supertype)}
Sometimes you keep completely separate relations for each subtype.
This can work when subtypes are operationally separate and rarely treated uniformly.
But it becomes awkward if you need cross-subtype reporting or shared relationships.
For teaching, this pattern is often less helpful as a default.

\section{Subtype modelling decision guide}
Deciding whether to use subtypes should be explicit.
Use the following criteria:
\begin{itemize}[leftmargin=*,nosep]
  \item subtype-specific attributes are substantial,
  \item subtype-specific constraints are important,
  \item subtype lifecycles differ (states or transitions),
  \item completeness/disjointness must be enforced,
  \item reports or rules treat subtypes differently.
\end{itemize}

\subsection{Attribute versus subtype}
If type is mostly for labeling or filtering, and there are no meaningful subtype-specific constraints, an attribute is often enough.
If correctness depends on subtype-specific rules or attributes, use subtypes.

\subsection{Optionality and completeness}
If subtype membership is incomplete, you allow a supertype instance with no subtype row.
If it is complete, you must ensure every supertype row has exactly one subtype row (disjoint) or at least one subtype row (overlapping).
Subtype structures are only correct when their disjointness and completeness rules are specified and validated.

\section{Worked example: Payment subtypes}
Assume payments can be made by card or bank transfer.
Each subtype has attributes that matter for correctness.

\begin{examplebox}
\textbf{Payment subtypes (complete, disjoint)}\par\smallskip
\[
\rel{Payment}(\underline{payment\_id}, amount, created\_at)
\]
\[
\rel{CardPayment}(\underline{payment\_id} \rightarrow \rel{Payment}, last4, auth\_code)
\]
\[
\rel{BankTransfer}(\underline{payment\_id} \rightarrow \rel{Payment}, iban, reference)
\]
Every payment is exactly one of these subtypes.
\end{examplebox}

\section{Completeness and disjointness test cases}
Subtype rules should be validated with simple tests.

Test 1: Can a \rel{Payment} exist without any subtype row? If not, completeness is required.\\
Test 2: Can a \rel{Payment} be both card and bank transfer? If not, disjointness is required.\\
Test 3: Are subtype-specific attributes allowed to be missing? If not, subtype tables are needed.

\section{Anti-patterns}
Subtype modelling has common traps.

EAV misuse: modelling subtypes as generic attribute/value rows hides constraints.\\
Type field with many nulls: a single table with dozens of optional columns weakens correctness.

\section{Common mistakes}
Mistakes here often lead to inconsistent data.

Mistake: Using a type attribute without enforcing subtype rules.\\
Fix: Add subtype tables or explicit constraints.

Mistake: Placing shared relationships on subtype tables.\\
Fix: Keep shared relationships on the supertype.

Mistake: Leaving disjointness/completeness unspecified.\\
Fix: State and validate these rules explicitly.

\section{Mini checklist}
Before finalizing subtypes, check:
\begin{itemize}[leftmargin=*,nosep]
  \item Are subtype-specific attributes meaningful and required?
  \item Is completeness specified (every supertype has a subtype)?
  \item Is disjointness specified (can instances be in multiple subtypes)?
  \item Are shared relationships placed on the supertype?
  \item Are subtype rules validated with test cases?
  \item Would a simple type attribute be sufficient instead?
  \item Are anti-patterns (EAV, many nulls) avoided?
\end{itemize}

\begin{chaptersummary}
Optionality is a statement about allowed states, not a nullable column choice.
Association relations represent many-to-many meaning and often become central entities when they carry attributes or time semantics.
Specialization (subtypes) is useful when “kinds of” entities have distinct attributes, constraints, or lifecycles.
Subtype models require explicit rules about disjointness and completeness, and the relational pattern should be chosen to preserve those rules clearly.
\end{chaptersummary}

\begin{exercisebox}
Answer in full sentences and then write the corresponding schema fragments in book notation.
\begin{enumerate}[label=\arabic*.]
  \item Give one example of a relationship where optionality is mandatory in your domain and justify why.
  \item Create a many-to-many relationship and design an association relation with at least two relationship attributes.
  \item Choose a concept (payment, appointment, user role) and decide whether to model it as a subtype or a type attribute. Justify your choice.
  \item If you chose subtypes, state whether they are disjoint and whether the specialization is complete.
  \item Design a supertype + subtype schema fragment using shared primary keys.
  \item Write two test cases that would detect violations of completeness or disjointness.
  \item Show an anti-pattern using a type field with many nulls and explain how you would fix it.
  \item Add subtype-specific constraints to your schema and explain where they belong.
\end{enumerate}
\end{exercisebox}

In the next chapter we formalize dependencies.
Subtypes and associations often reveal hidden functional dependencies.
